\chapter{Project Methodology}
\label{ch:methodology}

\begin{guidance}[What this chapter does]
Explains \emph{how} you ran the project: the approach for each stage, the
tools and methods you chose and why, and how you planned, monitored and
controlled the work. The Handbook asks for a flowchart showing the flow of
information across the main project stages, and for the Gantt chart and
risk analysis. Remember: the process is assessed as heavily as the product.
Typical length: 600--900 words.
\end{guidance}

\section{Methodology overview}
\label{sec:method-overview}

\Cref{fig:methodology} shows the stages of the project and the information
passed between them.

\begin{figure}[htbp]
  \centering
  \begin{tikzpicture}[node distance=7mm and 12mm]
    \node[fyp terminal] (start) {Project proposal};
    \node[fyp process, below=of start] (lit) {Literature review and market research};
    \node[fyp deliverable, below=of lit] (spec) {Product design specification};
    \node[fyp process, below=of spec] (concepts) {Generate candidate solutions};
    \node[fyp decision, below=of concepts] (select) {Meets specification?};
    \node[fyp process, right=of select] (build) {Implement prototype};
    \node[fyp process, above=of build] (test) {Test against requirements};
    \node[fyp decision, above=of test] (ok) {All tests passed?};
    \node[fyp terminal, above=of ok] (eval) {Evaluation and report};
    \node[fyp data, left=of lit] (log) {Logbook and project folder};

    \draw[fyp arrow] (start) -- (lit);
    \draw[fyp arrow] (lit) -- (spec);
    \draw[fyp arrow] (spec) -- (concepts);
    \draw[fyp arrow] (concepts) -- (select);
    \draw[fyp arrow] (select) -- node[above, font=\scriptsize]{yes} (build);
    \draw[fyp arrow] (select.west) -- ++(-6mm,0) |- node[pos=0.25, left, font=\scriptsize]{no} (concepts.west);
    \draw[fyp arrow] (build) -- (test);
    \draw[fyp arrow] (test) -- (ok);
    \draw[fyp arrow] (ok) -- node[right, font=\scriptsize]{yes} (eval);
    \draw[fyp arrow] (ok.east) -- ++(6mm,0) |- node[pos=0.25, right, font=\scriptsize]{no} (build.east);
    \draw[fyp arrow, dashed] (log) -- (lit);
  \end{tikzpicture}
  \caption{Project methodology and flow of information between stages (example)}
  \label{fig:methodology}
\end{figure}

\begin{guidance}
Describe each stage of the flowchart and justify the chosen methods, tools
and techniques with references or your own trials. Draw your own flowchart
with the styles \texttt{fyp process}, \texttt{fyp decision},
\texttt{fyp terminal}, \texttt{fyp io}, \texttt{fyp data},
\texttt{fyp deliverable} and \texttt{fyp arrow} defined in
\texttt{setup/fyp-common.tex}.
\end{guidance}

\section{Project plan}
\label{sec:plan}

\begin{guidance}
The Handbook asks for a Gantt chart that shows, as a minimum, a realistic
duration (in hours) and phasing of each task, the tasks grouped by project
stage (Proposal, Progress Report, Final Report, Engineering Project Day), and
the actual progress compared with the original plan. Use work breakdown
structure (WBS) numbers that match your logbook entries, so that the marker
can trace the plan to the evidence. A full-size chart can go in
\cref{app:project-management}. \Cref{fig:gantt-summary} shows planned
(grey) and actual (blue) durations; the planned hours add up to the 400
hours of the module (45 + 155 + 200).
\end{guidance}

\begin{figure}[htbp]
  \centering
  \begin{ganttchart}[fyp]{1}{30}
    \fypsemesters
    \ganttgroup{Proposal (45 h)}{1}{2}\\
    \ganttbar{1.1 Problem, aims, objectives (20 h)}{1}{2}
      \ganttbar[fyp actual]{}{1}{2}\\
    \ganttbar{1.2 Proposal, H\&S, ethics (25 h)}{1}{2}
      \ganttbar[fyp actual]{}{1}{3}\\
    \ganttmilestone{Project proposal}{2}\\
    \ganttgroup{Progress Report (155 h)}{3}{12}\\
    \ganttbar{2.1 Literature and market (50 h)}{3}{8}
      \ganttbar[fyp actual]{}{3}{9}\\
    \ganttbar{2.2 Specification (25 h)}{6}{9}
      \ganttbar[fyp actual]{}{7}{10}\\
    \ganttbar{2.3 Concepts, options (40 h)}{8}{11}
      \ganttbar[fyp actual]{}{9}{11}\\
    \ganttbar{2.4 Write Progress Report (40 h)}{10}{12}
      \ganttbar[fyp actual]{}{10}{12}\\
    \ganttmilestone{AE1 Progress Report}{12}\\
    \ganttgroup{Final Report (190 h)}{16}{27}\\
    \ganttbar{3.1 Implementation (80 h)}{16}{22}
      \ganttbar[fyp actual]{}{16}{24}\\
    \ganttbar{3.2 Testing (40 h)}{21}{24}
      \ganttbar[fyp actual]{}{23}{25}\\
    \ganttbar{3.3 Evaluation (20 h)}{24}{25}
      \ganttbar[fyp actual]{}{25}{26}\\
    \ganttbar{3.4 Write Final Report (50 h)}{23}{27}
      \ganttbar[fyp actual]{}{24}{27}\\
    \ganttmilestone{AE2 Final Report}{27}\\
    \ganttgroup{Engineering Project Day (10 h)}{28}{30}\\
    \ganttbar{4.1 Poster and demonstration (10 h)}{28}{30}
      \ganttbar[fyp actual]{}{28}{30}\\
    \ganttmilestone{AE3 Engineering Project Day}{30}
  \end{ganttchart}
  \caption{Project plan with WBS numbers and planned hours, grouped by project
    stage: planned (grey) and actual (blue) durations (example)}
  \label{fig:gantt-summary}
\end{figure}

\section{Risk management}
\label{sec:risk}

\begin{guidance}
Project risks (technical, resources, time, skills, equipment) in a risk
register: likelihood and severity on a scale of 1--5 (Very Unlikely to Very
Likely; Negligible to Catastrophic), risk rating = likelihood
$\times$ severity, and mitigation. The macro \verb|\riskrating{L}{S}|
calculates and colours the rating. Health and safety hazards are assessed
separately in a GRA01 form (\cref{app:hse}). Principles: BS ISO 31000
\parencite{bsi2018risk}.
\end{guidance}

\begin{table}[htbp]
  \centering
  \caption{Extract from the project risk register (example)}
  \label{tab:risk}
  \small
  \begin{tabular}{p{0.05\linewidth}p{0.3\linewidth}ccp{0.14\linewidth}p{0.25\linewidth}}
    \toprule
    ID & Risk & L & S & Rating & Mitigation \\
    \midrule
    1 & Sensor delivery delayed by more than two weeks & 3 & 3 & \riskrating{3}{3}
      & Order in week 4; identify an alternative supplier \\
    2 & Laboratory access limited during assessment weeks & 2 & 2 & \riskrating{2}{2}
      & Book sessions in advance \\
    3 & Prototype damaged during testing & 2 & 4 & \riskrating{2}{4}
      & Spare components; test at reduced load first \\
    4 & Final report underestimated & 4 & 4 & \riskrating{4}{4}
      & Start writing in Semester 2, week 6 \\
    \bottomrule
  \end{tabular}
\end{table}

\section{Ethics, health and safety}
\label{sec:ethics}

\begin{guidance}
Summarise the ethics release (Solent Ethics App), the DSE self-assessment and
any GRA01 risk assessments, and refer to the evidence in \cref{app:hse}.
Mention any ethical considerations, for example a user survey, data
protection or the Engineering Council Statement of Ethical Principles
\parencite{engc2026sep}.
\end{guidance}
